\documentclass[10pt,twocolumn]{article}

\usepackage[letterpaper,margin=0.72in,columnsep=0.27in]{geometry}
\usepackage[T1]{fontenc}
\usepackage{lmodern}
\usepackage{microtype}
\usepackage{amsmath,amssymb}
\usepackage{booktabs}
\usepackage{graphicx}
\usepackage{enumitem}
\usepackage{xcolor}
\usepackage{titlesec}
\usepackage[numbers,sort&compress]{natbib}
\usepackage[colorlinks=true,allcolors=blue!45!black]{hyperref}

\definecolor{ink}{HTML}{1F2933}
\definecolor{accent}{HTML}{294C60}
\color{ink}
\titleformat{\section}{\large\bfseries\color{accent}}{}{0pt}{}
\titleformat{\subsection}{\normalsize\bfseries\color{accent}}{}{0pt}{}
\titlespacing*{\section}{0pt}{0.8em}{0.35em}
\titlespacing*{\subsection}{0pt}{0.65em}{0.25em}
\setlength{\parindent}{1em}
\setlength{\parskip}{0.1em}
\setlist{nosep,leftmargin=1.25em}
\setlength{\bibsep}{1pt}
\renewcommand{\bibfont}{\footnotesize}

\newcommand{\LCDM}{$\Lambda$CDM}
\newcommand{\kmsmpc}{\mathrm{km\,s^{-1}\,Mpc^{-1}}}

\title{\vspace{-2.0em}\textbf{From Neighborhood Thinning to Cosmic Expansion:\\Quantitative Tests of a Relational-Complexity Hypothesis}}
\author{J. R. Landers}
\date{August 2026}

\begin{document}
\maketitle
\vspace{-1.2em}

\begin{abstract}
Adding genuinely distinguishing coordinates to an object increases its
separation from alternatives and thins every fixed-radius neighborhood under a
broad class of additive distances.  The result invites a speculative but
testable question: could physical expansion reflect the growth of relational
distinguishability?  We construct the simplest bridge by identifying inverse
exact-match density with relational volume, obtaining
$H=(\ln 2/3)\dot N$.  A minimal history law then predicts
$H_0=70.85\,\kmsmpc$ from the measured cosmic age, placing the observed Hubble
scale within a few percent without inserting $H_0$.  The coincidence is
instructive but incomplete: the same law produces a coasting universe and is
strongly rejected by the DESI DR2 baryon-acoustic-oscillation distance shape.
A two-channel extension fits the BAO data well, but is algebraically identical
to the flat-\LCDM{} background.  Replacing accumulated time by relativistic
causal-past four-volume also fails to yield a constant expansion--complexity
coupling.  Finally, public CAMELS simulations show that neighborhood change is
enhanced in dense, bound environments where expansion is suppressed.  Most of
this signal is predicted by peculiar motion, and it persists in matched
gravity-only simulations.  These results do not eliminate a relational account
of spacetime.  They narrow its viable content: the tested complexity measures
behave more naturally as diagnostics of gravitational dependency and binding
than as a local source of dark-energy expansion.
\end{abstract}

\section{A geometric question with cosmological stakes}

The discovery that the cosmic expansion is accelerating transformed the
cosmological constant from a formal possibility into an empirical problem
\citep{Riess1998,Perlmutter1999}.  Flat \LCDM{} remains an economical and highly
successful description of the background, particularly when confronted with
the cosmic microwave background and baryon acoustic oscillations
\citep{Planck2018,DESI2025}.  Yet its dark-energy scale is entered as a fitted
constant rather than derived from a microscopic principle.  That gap has kept
alive a broad question: might geometry be an effective description of deeper
statistical, thermodynamic, or informational structure?

There is serious precedent for asking the question carefully.  Einstein's
equation can be recovered as a local thermodynamic equation of state under
specific horizon-entropy assumptions \citep{Jacobson1995}, while causal-set
theory explores how Lorentzian order and spacetime volume might emerge from a
locally finite relational structure \citep{Surya2019}.  These programs do not
license an arbitrary identification of information with geometry.  They do,
however, show what a successful bridge must eventually provide: a covariant
object, a dynamical law, and quantitative contact with observation.

The present work begins from a much simpler exact fact.  Under additive
nonnegative distances, appending a distinguishing coordinate can only increase
pairwise separation and can only reduce the population inside a fixed-radius
neighborhood.  For independent binary coordinates, exact-match density falls
as $2^{-N}$.  If the inverse of this density is interpreted as an effective
relational volume, then complexity growth has the kinematics of expansion.

Our purpose is not to promote that analogy by rhetoric.  It is to give it every
reasonable quantitative opportunity to succeed, then identify precisely what
survives.  We proceed in four stages: derive the exact kinematic bridge, test
simple history laws against the Hubble scale and the DESI DR2 distance shape,
replace naive accumulated time by a relativistic causal-history construction,
and examine neighborhood evolution directly in hydrodynamic and gravity-only
cosmological simulations.

\section{The exact bridge---and the missing dynamics}

Consider objects described by $N$ coordinates, with additive distance
\begin{equation}
 \delta_N(O_i,O_j)=\sum_{n=1}^{N}d_n(o_{in},o_{jn}),
 \qquad d_n\geq0.
\end{equation}
Appending a coordinate gives
$\delta_{N+1}=\delta_N+d_{N+1}\geq\delta_N$.  Consequently every
fixed-radius neighborhood is nonincreasing with $N$.  In the particularly
transparent case of independent uniform binary coordinates and exact matching,
the expected neighborhood density is
\begin{equation}
 \rho_N=2^{-N}.
\end{equation}
Define an effective relational volume and a three-dimensional scale factor,
\begin{equation}
 V_{\rm rel}=\rho_N^{-1},\qquad
 a_{\rm rel}=V_{\rm rel}^{1/3}=2^{N/3}.
\end{equation}
The corresponding expansion rate is then exact:
\begin{equation}
 H_{\rm rel}\equiv\frac{\dot a_{\rm rel}}{a_{\rm rel}}
 =\frac{\ln2}{3}\dot N.
 \label{eq:bridge}
\end{equation}

Equation~\eqref{eq:bridge} is the central bridge, but also the central warning.
It is a dictionary, not a field equation.  The factor of three comes from the
chosen volume interpretation; the bit rate $\dot N$ has not been derived; and
no stress--energy tensor or conservation law has appeared.  Physical content
begins only when $N$ is defined independently of the observed scale factor and
its dynamics predict more than one number.

\section{A numerical coincidence worth testing}

Suppose that maintaining a distinguishable three-dimensional history requires
one new binary distinction per spatial dimension whenever the temporal
baseline doubles:
\begin{equation}
 N(t)=3\log_2(t/t_*).
 \label{eq:agelaw}
\end{equation}
Equations~\eqref{eq:bridge} and \eqref{eq:agelaw} imply
\begin{equation}
 a(t)\propto t,\qquad H(t)=t^{-1}.
 \label{eq:coasting}
\end{equation}
Using $t_0=13.80\,\mathrm{Gyr}$ gives
\begin{equation}
 H_0^{\rm age}=70.85\,\kmsmpc.
\end{equation}
This is $5.1\%$ above the Planck base-\LCDM{} value
$67.4\pm0.5\,\kmsmpc$ \citep{Planck2018}, $3.9\%$ above the DESI
DR2+CMB value used here, and $3.0\%$ below the SH0ES distance-ladder value
$73.04\pm1.04\,\kmsmpc$ \citep{Riess2022}.  Because CMB-derived age and
$H_0$ are correlated and model dependent, this is not a precision independent
prediction.  It is nevertheless a striking dimensional result: a transparent
information law lands on the observed Hubble scale without inserting it.

The same simplicity makes the decisive failure immediate.  Equation
\eqref{eq:coasting} gives $\ddot a=0$ and $q=0$, whereas the measured late
universe accelerates.  A correct value of $H$ at one epoch is not a correct
expansion history.  The coincidence therefore earns a shape test rather than a
claim.

\section{DESI tests the history, not merely the scale}

Baryon acoustic oscillations provide a standard-ruler measurement of transverse
and radial cosmic distances \citep{Eisenstein2005}.  We use the published
13-element DESI DR2 BAO data vector and covariance \citep{DESI2025}.  For every
model, the nuisance normalization
\begin{equation}
 X=\frac{c}{H_0r_d}
\end{equation}
is profiled analytically.  The comparison therefore tests the redshift
dependence of the distance relation without using the value of $H_0$ that
motivated the hypothesis.

For the family
\begin{equation}
 N(t)=d_{\rm eff}\log_2(t/t_*),\quad
 a(t)\propto t^p,\quad p=d_{\rm eff}/3,
\end{equation}
we test a matter-like history, the coasting spatial history, a literal
``3D plus time'' history, and a freely fitted constant $p$.  Table~\ref{tab:bao}
and Fig.~\ref{fig:desi} show the result.  No constant-dimensional history
describes the BAO shape.  Even the best constant power gives
$\chi^2=96.97$ for 11 degrees of freedom.  More sharply, fitting $p$ and $X$
only below $z=1$ gives $p=1.1185$, but the six held-out measurements at
$z\geq1$ yield $\chi^2=670.72$ without refitting.

\begin{table}[t]
\centering
\caption{DESI DR2 BAO shape fits.  The BAO normalization is fitted in every
row.  The two-channel row also fits its channel ratio.}
\label{tab:bao}
\small
\begin{tabular}{lrr}
\toprule
Model & $\chi^2$ & dof\\
\midrule
$p=2/3$ & 1420.18 & 12\\
$p=1$ & 187.09 & 12\\
$p=4/3$ & 1364.93 & 12\\
Best constant $p=0.9125$ & 96.97 & 11\\
Two channel & 10.27 & 11\\
Flat \LCDM{}, $\Omega_m=0.3027$ & 10.63 & 12\\
\bottomrule
\end{tabular}
\end{table}

\begin{figure*}[t]
\centering
\includegraphics[width=0.82\textwidth]{figures/desi_model_comparison.pdf}
\caption{Reduced $\chi^2$ for the DESI DR2 BAO distance shape after profiling
the common ruler normalization.  The logarithmic axis makes visible both the
strong rejection of constant-dimensional histories and the good background fit
of the two-channel and flat-\LCDM{} forms.}
\label{fig:desi}
\end{figure*}

\subsection{The successful extension and its honest interpretation}

A more flexible relational construction separates local material interactions
from a persistent history channel and combines statistically independent rates
in quadrature:
\begin{equation}
 \dot N_{\rm eff}^{,2}
 =\dot N_{m,0}^{,2}a^{-3}+\dot N_h^{,2}.
 \label{eq:channels}
\end{equation}
Together with Eq.~\eqref{eq:bridge},
\begin{equation}
 H^2(a)=\left(\frac{\ln2}{3}\right)^2
 \left(\dot N_{m,0}^{,2}a^{-3}+\dot N_h^{,2}\right).
 \label{eq:twochannel}
\end{equation}
Fitting the channel ratio and BAO scale gives
$\Omega_{m,\rm eff}=0.2975$ and $\chi^2=10.27$ for 11 degrees of freedom.
A fit restricted to $z<1$ gives $\Omega_{m,\rm eff}=0.3095$ and predicts the
six high-redshift measurements with $\chi^2=4.74$.  With the reference ruler
$r_d=147.09\,\mathrm{Mpc}$, the full fit corresponds to
$H_0=69.03\,\kmsmpc$, a present material-channel rate of
$0.1667\,\mathrm{bits\,Gyr^{-1}}$, and a history-channel rate of
$0.2561\,\mathrm{bits\,Gyr^{-1}}$.

This is a useful success of form, but not yet a new explanation.  Equation
\eqref{eq:twochannel} is exactly the flat-\LCDM{} background equation under a
change of names.  It predicts the familiar acceleration transition,
$z_{\rm acc}\simeq0.68$, but becomes explanatory only if the quadratic
composition rule, the $a^{-3}$ term, the constant rate, and the approximately
$0.30{:}0.70$ channel ratio follow from an independent microscopic theory.
The fit tells us what must be derived; it does not perform the derivation.

\section{Time must be made relativistic}

The literal addition of time as a fourth positive metadata coordinate is not
covariant.  Cosmic time is not an interchangeable Euclidean feature, and
relativistic observables are organized by causal structure.  Motivated by the
role of causal order in discrete approaches to spacetime \citep{Surya2019}, we
replace accumulated coordinate time by the causal past of an observation event.
In flat FLRW spacetime its four-volume is
\begin{equation}
 V_4(\tau_o)=\frac{4\pi}{3}\int_0^{\tau_o}\!d\tau'\,a^3(\tau')
 \left[\int_{\tau'}^{\tau_o}\frac{c\,d\tau''}{a(\tau'')}\right]^3.
 \label{eq:v4}
\end{equation}
The covariant expansion scalar is $\theta=\nabla_\mu u^\mu=3H$.  We ask
whether a redshift-independent constitutive relation
\begin{equation}
 3H=\kappa\dot C
 \label{eq:kappa}
\end{equation}
can hold for either $C=V_4$ in fixed units or $C=\log_2V_4$.  Evaluating
Eq.~\eqref{eq:v4} from $z=2.33$ to the present along the DESI-calibrated
two-channel background, the required $\kappa$ varies by a factor of 321.8 for
raw four-volume.  Logarithmic volume performs substantially better, but still
requires a factor 1.514 drift.  Its present coupling is $0.5373$, rather than
the $\ln2=0.6931$ supplied by the direct neighborhood-volume dictionary.

This result is more constructive than simply declaring ``time is different.''
It identifies the problem: total causal volume repeatedly counts shared and
redundant history.  If a history channel exists, the more plausible object is
the \emph{conditional new information} between neighboring causal slices after
conditioning on their common past.  That object is not supplied by
Eq.~\eqref{eq:v4}, so we next seek an empirical analogue in evolving
cosmological neighborhoods.

\section{Neighborhood evolution in CAMELS}

CAMELS provides many $25\,h^{-1}\mathrm{Mpc}$ cosmological volumes, including
IllustrisTNG hydrodynamic simulations and gravity-only counterparts
\citep{CAMELS2021,CAMELSRelease2022}.  The hydrodynamic runs use the
IllustrisTNG galaxy-formation framework \citep{Pillepich2018}.  These
simulations assume an FLRW/\LCDM{} background and therefore cannot establish
that information causes cosmic expansion.  They can test the narrower local
prediction: does independently measured neighborhood novelty rise where
physical expansion is stronger, and is it suppressed in bound regions?

\subsection{Tracking and estimators}

We analyze snapshots 86, 88, and 90, corresponding approximately to
$z=0.100$, $0.049$, and 0.  Subhalos with at least 100 particles are matched
between adjacent snapshots using mutual-nearest comoving positions, a mass
ratio within a factor of three, and a maximum displacement of
$0.5\,\mathrm{Mpc}/h$.  Where the most-bound particle remains available, this
geometric procedure recovers the persistent identity in 94--98\% of cases
across the analyzed samples.

For every matched object we construct a $k=16$ nearest-neighbor set.  A discrete
estimator assigns surprisal to newly appearing neighbors conditional on the old
set.  A continuous robustness estimator assigns Gaussian weights over the
union of 64 nearby candidates and measures the Jensen--Shannon divergence
between successive neighborhood distributions \citep{Lin1991}.  The physical
volume change follows the same old neighbors into the next snapshot, using the
determinant of their local covariance ellipsoid and restoring the background
$a^3$ factor.  Low-density centrals define the field sample.  Satellites or
objects in the upper local-density quartile define a bound-region proxy.

In the initial fixed-cosmology volume, neighbor retention is 96.5--96.6\%.
The field follows global expansion, with local-to-global ratios 0.982 and
0.989, while the bound proxy gives 0.460 and 0.455.  Yet the continuous
neighborhood JS change is roughly 24 times larger in the bound proxy.  Novelty
correlates positively with density and negatively with local expansion.  The
same environmental separation persists for $k=8,32,$ and 64.  Thus the direct
prediction has the wrong local sign: the largest relational rearrangements
occur not in freely expanding neighborhoods but in structured regions of
orbital motion, close passages, and mergers.

\subsection{Conditioning on phase space}

Much of this rearrangement can be reversible dynamics rather than new
information.  We therefore propagate each old phase-space state using its
measured peculiar velocity.  In comoving $\mathrm{Mpc}/h$, the first-order FLRW
prediction is
\begin{equation}
 \begin{split}
 \boldsymbol{x}_{\rm pred}(a_2)&=\boldsymbol{x}(a_1)
 +\frac{\boldsymbol{v}_{\rm pec}(a_1)}{100}
 \int_{a_1}^{a_2}\frac{da}{a^2E(a)},\\
 E(a)&=\sqrt{\Omega_ma^{-3}+\Omega_\Lambda}.
 \end{split}
 \label{eq:phase}
\end{equation}
The residual is the JS divergence between the neighborhood predicted by
Eq.~\eqref{eq:phase} and the observed neighborhood.  This is a deliberately
conservative ballistic control: accelerations, mergers, disruption, and halo
identification remain in the residual.

Across five independent fixed-cosmology IllustrisTNG realizations and two
intervals per realization, field patches expand at $0.992\pm0.021$ of the
global rate, whereas the bound proxy expands at $0.500\pm0.089$.  Raw median JS
changes are $(3.80\pm0.85)\times10^{-5}$ bits in the field and
$(8.39\pm1.53)\times10^{-4}$ bits in the bound proxy.  Phase conditioning
reduces these to $(1.33\pm0.32)\times10^{-6}$ and
$(1.40\pm0.64)\times10^{-4}$ bits: reductions of 96.5\% and 83.3\%,
respectively.  The residual remains anticorrelated with expansion
($\rho=-0.304\pm0.047$) and correlated with density
($\rho=+0.629\pm0.032$).  Quoted dispersions are sample standard deviations
across realization--interval measurements, not cosmological parameter errors.

\begin{figure*}[t]
\centering
\includegraphics[width=0.82\textwidth]{figures/camels_phase_control.pdf}
\caption{Continuous neighborhood change in field and bound-proxy samples.
Velocity propagation removes most of the raw change, especially in the field.
The bound excess and its survival in gravity-only simulations indicate that
ordinary nonlinear gravitational dynamics, not baryonic complexity, is the
dominant source of the measured signal.  Error bars show sample standard
deviations across realization--interval measurements.}
\label{fig:camels}
\end{figure*}

\subsection{A matched gravity-only control}

To isolate baryonic structural complexity, we repeat the analysis for matched
IllustrisTNG and IllustrisTNG-DM Latin-hypercube realizations LH\_0--LH\_2,
which share simulation indices, cosmologies, and initial seeds.  The
hydrodynamic and gravity-only samples yield field expansion ratios
$1.000\pm0.011$ and $0.988\pm0.017$, and bound ratios $0.533\pm0.062$ and
$0.542\pm0.049$.  Their raw field JS changes are nearly identical.  The raw
bound change is $(1.041\pm0.046)\times10^{-3}$ bits in hydrodynamics and
$(1.241\pm0.116)\times10^{-3}$ bits in gravity only; after phase conditioning
the corresponding residuals are $(2.32\pm1.34)\times10^{-4}$ and
$(3.21\pm0.56)\times10^{-4}$ bits (Fig.~\ref{fig:camels}).

The qualitative effect therefore does not require gas physics, star formation,
or feedback.  Collisionless gravity is sufficient.  This does not make the
relational pattern uninteresting.  It changes its most natural interpretation:
neighborhood complexity records the dynamically maintained dependencies of
bound structure rather than acting as an independent pressure that expands
space.

\section{What the tests leave standing}

The original intuition separates into three claims that should no longer be
conflated.

First, \emph{structural distinction thins abstract neighborhoods}.  This is an
exact combinatorial statement under the stated distance assumptions, and none
of the cosmological tests alters it.

Second, \emph{inverse neighborhood density can be written as an effective
volume}.  This produces the exact kinematic relation
Eq.~\eqref{eq:bridge}.  It is mathematically coherent and explains why the
Hubble scale emerges so readily from a bit rate.  But without independent
dynamics, any measured $H(t)$ can be translated into a fitted $N(t)$, making
the construction tautological.

Third, \emph{complexity growth physically drives cosmic expansion}.  The
specific realizations tested here do not support this claim.  Fixed-dimensional
history laws fail the BAO shape; total causal history lacks a constant
coupling; and direct local novelty is largest where expansion is smallest.  A
successful future model must introduce a covariant, coarse-grained,
irreversible information functional and derive its gravitational coupling
before using expansion data.  It must also distinguish itself from \LCDM{} in
growth, lensing, redshift-space distortions, or another observable, rather than
only relabeling the background equation.

The binding result offers a productive direction.  In standard gravity, the
competition between a bound mass $M$ and late-time accelerated expansion sets
a maximum turnaround scale
\begin{equation}
 R_{\rm TA,max}=\left(\frac{3GM}{\Lambda c^2}\right)^{1/3},
\end{equation}
which has been proposed as a test of \LCDM{} and alternatives
\citep{Pavlidou2014}.  The CAMELS results suggest that relational dependency
measures may illuminate the structured side of this boundary.  The challenge
is to derive---not rename---the dependence on $G$, $M$, and $R$, and to show
whether an information balance reproduces the turnaround normalization.

\section{Conclusion}

The neighborhood-thinning idea survives contact with cosmology, but in a more
interesting and disciplined form than the opening analogy.  Its exact
combinatorics provide a clean bridge from distinguishing information to an
effective expansion rate.  A minimal history law then lands unexpectedly close
to the observed Hubble scale.  Precision distance data show that the same law
cannot describe the expansion history, while a successful two-channel version
reveals itself as the standard flat-\LCDM{} background in relational language.
Making history relativistic improves the formulation but not yet the dynamics.

The simulations supply the most concrete new clue.  Relational rearrangement
is enhanced where gravity binds and cosmic expansion is suppressed; ordinary
phase-space motion explains most of it; and gravity alone reproduces it.  The
result redirects the hypothesis rather than ending it.  Neighborhood
complexity appears better suited to describing how gravitational structures
maintain dependent histories than to serving, in its present form, as the
source of dark energy.  That narrowed target is now precise: construct an
irreversible covariant information measure, derive its coupling, and demand a
prediction that cannot be obtained by translating the observed Friedmann
equation.

\section*{Data and reproducibility}

All calculations are implemented in the accompanying repository.  The DESI
test uses the public DR2 BAO mean vector and covariance.  CAMELS group catalogs
are downloaded from the public release \citep{CAMELSRelease2022}; positions and
velocities follow the published Illustris catalog conventions.  The scripts
record the selections, matching cuts, estimators, phase propagation, and all
reported fit statistics.  No CAMELS result is used as an independent test of
the assumed background cosmology.

\bibliographystyle{unsrtnat}
\bibliography{complexity_cosmology}

\end{document}
